\section{总结与延伸}

\subsection{一条贯穿全课的系统主线}

这堂 overview 可以压缩成一条依赖链：原始世界先被 tokenization/encoder 映射为表示；architecture 决定信息怎样交互；pretraining data 与 objective 决定参数中形成什么统计规律；post-training 与 inference-time compute 改变行为；retrieval、memory、tools 和 agent loop 增加外部状态；interpretability、alignment 与 evidence protocol 决定系统能否被理解和约束；JEPA 与 SSM 则分别探索新的预测对象和计算路径。

这条链也解释了为什么系统故障经常跨层传播。Tokenizer 破坏实体边界，会让 retrieval query 变差；检索返回错误证据，会让 reasoning trace 在错误前提上变得更长；preference model 若偏爱自信语气，agent 可能减少必要的 abstention；memory 若保存未经验证的总结，后续每次任务都会把一次错误变成持久状态。相反，局部修复也可能失效：只扩大 context 不能修复错误权限，只做 DPO 不能更新过期知识，只加 verifier 不能补回不可观察变量。工程审计应沿 representation、data、objective、state、action 与 evaluation 逐层定位，而不是把所有问题归为“模型不够聪明”。

因此，一份可信系统报告至少要同时给出训练数据时间边界、模型与 tokenizer 版本、post-training 方法、可用工具与权限、memory 生命周期、评估集和成本指标。缺少这些信息时，即使复现了单个 benchmark 分数，也无法复现真实系统行为；模型名称不是完整实验条件。

同样，任何“替代 Transformer”的主张都应说明比较发生在何种序列长度、batch、硬件、训练预算与质量阈值下。离开 workload 的架构排名没有稳定含义，真正可迁移的是测量方法与边界条件。

\begin{importantbox}{全课最可迁移的十问}
\begin{enumerate}
  \item 输入被切成什么 token，哪些信息在入口就丢失？
  \item 目标函数奖励什么统计规律，遗漏什么？
  \item 数据来源、混合、顺序与更新频率是什么？
  \item 计算瓶颈在参数、attention、memory、retrieval 还是工具？
  \item Post-training 改权重、改候选搜索，还是加外部状态？
  \item Agent 的 observation、action、权限、budget 与停止条件是什么？
  \item 指标测的是 loss、任务成功、偏好还是现实影响？
  \item 错误来自 belief、observation、planning 还是 reward proxy？
  \item 系统如何保存、更新、冲突解决和遗忘？
  \item 新架构在同等数据、compute 与 serving 条件下改善了哪个 Pareto 维度？
\end{enumerate}
\end{importantbox}

\subsection{三条证据边界}

第一，课堂中的 BabyLM、RAG、CGLS、fMRI 与 hallucination 研究都对应具体实验设计，不能从单一规模或数据集外推为普遍规律。第二，应用 demo 与 benchmark 不能替代真实工作流中的成本、权限、分布偏移和外部验证。第三，world model、continual learning、alignment 等概念必须先定义更新对象和评价协议，否则同一词会覆盖互不相同的问题。

\subsection{开放问题}

下一步值得继续追问：能否在不牺牲可随机访问能力的前提下获得线性长序列计算？能否把 retrieval、parameter update 与 memory conflict 统一成可验证的持续学习系统？能否构造对 world-model belief、planner 与 action 分层诊断的 agent benchmark？能否让 process supervision 与 mechanistic evidence 相互校验，而不是共同奖励表面 reasoning style？

\subsection{拓展阅读}

\begin{itemize}
  \item 官方课程：\url{https://web.stanford.edu/class/cs25/}
  \item 官方录像：\url{https://www.youtube.com/watch?v=bHSDPgZYie0}
  \item Baby Scale：\url{https://arxiv.org/abs/2603.29522}
  \item Bilingual BabyLM：\url{https://arxiv.org/abs/2603.29552}
  \item Pretraining--RAG scaling：\url{https://arxiv.org/abs/2604.00715}
  \item Curriculum-Guided Layer Scaling：\url{https://arxiv.org/abs/2506.11389}
  \item Cross-network fMRI representation：\url{https://arxiv.org/abs/2603.00786}
  \item Unified hallucination definition：\url{https://arxiv.org/abs/2512.21577}
\end{itemize}

\end{document}
